\subsection{Modul RTL \texttt{i2s\_filter\_bridge}}

IP \path{i2s_0} menerima sampel PCM melalui register AXI4-Lite dan menghasilkan aliran serial I2S, sedangkan modul \path{filter_coeff} memakai antarmuka paralel 16 bit berupa \path{x_in}, \path{valid}, dan \path{y_out}. Kedua antarmuka tersebut tidak dapat dihubungkan dengan sebuah net pada block design. Penelitian ini menambahkan modul \path{i2s_filter_bridge} pada keluaran serial IP I2S tanpa mengubah RTL IP I2S maupun RTL filter.

Pada tepi naik BCLK, bridge mendeteksi transisi WS, membuang satu bit tunda sesuai format Philips I2S, lalu menggeser 16 bit payload ke register sampel. Bridge membangkitkan \path{left_valid} atau \path{right_valid} setelah menerima satu word kanal. Block design memakai dua instans \path{filter_coeff} agar kanal kiri dan kanan memiliki state biquad yang terpisah.

Bridge menyimpan keluaran filter dan mengirimkannya pada frame berikutnya melalui \path{i2s_data_out}. Proses serialisasi memperbarui SDATA pada tepi turun BCLK agar DAC membaca nilai yang stabil pada tepi naik. Jalur MCLK, BCLK, dan WS menggunakan koneksi tembus; bridge hanya mengganti SDATA dengan data hasil filter. Konsekuensinya, integrasi menambah latensi satu frame stereo tanpa mengubah frekuensi sampling atau relasi fase antarsinyal I2S.

Tiga switch pemilih koefisien berasal dari pin FPGA yang asinkron terhadap clock sistem. Dua tingkat flip-flop pada bridge menyinkronkan nilai switch sebelum kedua filter membacanya. Listing~\ref{lst:ch5-i2s-filter-bridge} memperlihatkan antarmuka dan bagian utama jalur data bridge.

\begin{lstlisting}[
	language=Verilog,
	caption={Bagian utama RTL \texttt{i2s\_filter\_bridge}.},
	label={lst:ch5-i2s-filter-bridge}
]
module i2s_filter_bridge (
    input  wire               clk,
    input  wire               resetn,
    input  wire [2:0]         filter_sw_in,
    input  wire               i2s_bclk_in,
    input  wire               i2s_ws_in,
    input  wire               i2s_data_in,
    output reg                i2s_data_out,
    output reg                left_valid,
    output reg signed [15:0]  left_sample,
    input  wire signed [15:0] left_filtered,
    output reg                right_valid,
    output reg signed [15:0]  right_sample,
    input  wire signed [15:0] right_filtered
);

wire bclk_rise =  i2s_bclk_in & ~bclk_d;
wire bclk_fall = ~i2s_bclk_in &  bclk_d;

always @(posedge clk) begin
    if (!resetn) begin
        left_valid  <= 1'b0;
        right_valid <= 1'b0;
        i2s_data_out <= 1'b0;
    end else begin
        bclk_d      <= i2s_bclk_in;
        left_valid  <= 1'b0;
        right_valid <= 1'b0;

        /* Deserialisasi: abaikan bit tunda, ambil 16 bit payload. */
        if (bclk_rise) begin
            if (i2s_ws_in != rx_ws) begin
                rx_ws      <= i2s_ws_in;
                rx_bit_pos <= 6'd0;
            end else if (rx_bit_pos < 6'd16) begin
                rx_shift   <= {rx_shift[14:0], i2s_data_in};
                rx_bit_pos <= rx_bit_pos + 6'd1;
                if (rx_bit_pos == 6'd15) begin
                    if (!i2s_ws_in) begin
                        left_sample <= {rx_shift[14:0], i2s_data_in};
                        left_valid  <= 1'b1;
                    end else begin
                        right_sample <= {rx_shift[14:0], i2s_data_in};
                        right_valid  <= 1'b1;
                    end
                end
            end
        end

        /* Serialisasi hasil filter pada tepi turun BCLK. */
        if (bclk_fall) begin
            if (i2s_ws_in != tx_ws) begin
                tx_ws        <= i2s_ws_in;
                tx_bit_pos   <= 6'd0;
                i2s_data_out <= 1'b0;
            end else if (tx_bit_pos < 6'd16) begin
                i2s_data_out <= i2s_ws_in
                              ? tx_right[15-tx_bit_pos]
                              : tx_left[15-tx_bit_pos];
                tx_bit_pos <= tx_bit_pos + 6'd1;
            end else begin
                i2s_data_out <= 1'b0;
            end
        end
    end
end
endmodule
\end{lstlisting}

Pengujian mengirim nada sinus 1~kHz dari firmware ke register data IP I2S. Capture timing menunjukkan pulsa \path{left_valid} dan \path{right_valid} muncul satu kali per kanal pada setiap frame, sedangkan BCLK dan WS pada sisi DAC tetap mengikuti keluaran IP I2S. Pengamatan keluaran DAC mengonfirmasi bahwa bridge mempertahankan batas kanal dan urutan bit selama filter memproses aliran audio.
